\section{Certified computation from rational witnesses}
\label{sec:a-cert}

The additive algorithms above fix a tolerance before computation. A
numerical solver raises a different question: what can be proved about
its actual output? This section gives rational witnesses for energy,
individual flows, linear flow goals, and recovered resistance scenarios.
They are verified in exact arithmetic, without assuming that a numerical
method has reached its stopping tolerance.
The global resistance-design certificate in
\cref{prop:a-corr-certificate} already treats maximum dissipation over a
polytope; the present certificates concern the physical state of a fixed
law and its use in the arc-envelope reduction.

The underlying principles are classical. Conjugates and Slater duality
are used in their standard sense
\citep[Sections 3.3, 5.1.6, and 5.2.3]{boyd2004-convex}.
The divergence below is precisely the convex-function divergence of
\citet[p.~206, equation (1.4)]{bregman1967-relaxation}, with its second
argument equal to the unknown physical state. Primal-dual gaps as
computable bounds on uniformly convex errors are established a posteriori
methods; see \citet[Section 3, Proposition 3.1]{bartels2020-primal-dual}.
The conservation projection is the electrical cut/cycle projection
\citep[Chapter 2, Section 4]{lyons2016-probability}.
Our purpose is to give their complete asymmetric-quadratic specialization,
exact rational checks, degeneracy conditions, and original-instance
interpretation. We make no novelty claim for these general principles.

The accompanying Lean formalization (topics 03 and 16) proves the
deterministic energy, Bregman, support, and
curvature guarantees below, including conditional scenario recovery and
the two-path comparison. They do not cover the original
uncertainty-to-envelope mapping, the numerical producers, the Python
verifiers and parsers, conservation repair, the bit-complexity claims,
or the effective-resistance reformulation. In the formalization, the
rational goal-acceptance predicate checks a supplied factor $C(v)$;
optimality of $C_*$ is proved by separate theorems.

\subsection{Energy witnesses and rational construction}

Let $A$ have $+1$ at the tail and $-1$ at the head of each oriented
edge. In this section the graph may be a finite loopless multigraph,
possibly disconnected, with $m\ge1$; nominations must be feasible
on each component. Define
\begin{equation}
 \begin{split}
 E(x)&=\sum_e E_e(x_e),\qquad
 E_e(t)=\frac13
 \begin{cases}c_e^+t^3&t\ge0,\\c_e^-(-t)^3&t<0,\end{cases}\\
 G_e(t)&=E_e'(t)=c_e^{\sgn(t)}t|t|,\qquad
 \beta_L=\min_{e,\sigma\in\{+,-\}}c_e^\sigma>0.
 \end{split}
 \label{eq:a-cert-energy}
\end{equation}
All coefficients and nominations are rational. The notation
$c_e^{\sgn(t)}$ selects $c_e^+$ at $t=0$; every formula at zero
is independent of this convention. Strict convexity and coercivity
imply that $E$ has a unique minimizer $x^*$ on $Ax=b$.
Stationarity on this affine space gives
$G(x^*)=A^\top\pi^*$, with one arbitrary potential gauge per component.
Thus $x^*$ is the passive physical flow.

\begin{proposition}[Rational energy and uniform flow certificate]
\label{prop:a-cert-energy}
Let $y\in\Q^m$ satisfy $Ay=b$ and let $p\in\Q^V$ be arbitrary.
For $d=A^\top p$, suppose rational $\zeta_e\ge0$ satisfy
\[
 \zeta_e^2 c_e^{\sgn(d_e)}\ge |d_e|^3.
\]
Set
\begin{equation}
 L=b^\top p-\frac23\sum_e \zeta_e,\qquad
 \delta=E(y)-L.
 \label{eq:a-cert-gap}
\end{equation}
Then $0\le E(y)-E(x^*)\le\delta$. Any rational $\eta\ge0$
with $\beta_L\eta^3\ge6\delta$ gives
$|y_e-x_e^*|\le\eta$ for every edge. All hypotheses, the claimed gap,
and the radius condition are verifiable by rational arithmetic in
polynomial time.
\end{proposition}
\begin{proof}
For $d>0$ the maximizer of $dt-E_e(t)$ is positive, since negative
$t$ gives a nonpositive value whereas small positive $t$ gives a
positive value. Differentiation gives $t=\sqrt{d/c_e^+}$ and value
$(2/3)\sqrt{d^3/c_e^+}$. The negative case uses $c_e^-$ and the
maximizer $-\sqrt{|d|/c_e^-}$; at zero the maximizer and value
are zero. Hence
\[
 E_e^*(d)=\frac23\sqrt{\frac{|d|^3}{c_e^{\sgn(d)}}}.
\]
The defining inequality $E_e(t)+E_e^*(d)\ge dt$, summed over the
physical flow, proves $E(x^*)\ge L$. Minimality gives
$E(y)\ge E(x^*)$.

The scalar modulus used in \cref{eq:a-bound-energy-error} also holds
for all laws \eqref{eq:a-cert-energy}. The function
$G_e(t)-\beta_Lt|t|$ is nondecreasing, and
\[
 (u|u|-v|v|)(u-v)\ge\frac12|u-v|^3.
\]
On either same-sign half-line this follows by factoring; for opposite
signs it is $|u|^2+|v|^2\ge(|u|+|v|)^2/2$.
Writing $a=y_e-z_e$ and integrating with $u=z_e+ta$, $v=z_e$,
$0\le t\le1$, yields
\begin{equation}
 D_e(y_e,z_e):=E_e(y_e)-E_e(z_e)-G_e(z_e)(y_e-z_e)
 \ge\frac{\beta_L}{6}|y_e-z_e|^3.
 \label{eq:a-cert-scalar-modulus}
\end{equation}
Indeed the integrand $[G_e(z_e+ta)-G_e(z_e)]a$ is at least
$\beta_Lt^2|a|^3/2$. At $z=x^*$ the linear terms sum to zero
by $A(y-x^*)=0$. Summation proves the asserted radius.
The checks require only rational additions, products, comparisons, and
cubes, with polynomially bounded intermediate bit lengths.
\end{proof}

Here and below, verification complexity counts the expanded binary
lengths of rational numerators and denominators. A parser may accept
other exact string forms, including scientific exponent notation; the
bound applies after their expansion, not to their raw JSON length.
Neither $y$ nor $p$ needs to satisfy the constitutive equations; an
inaccurate candidate simply has a larger certified gap.

A practical construction repairs conservation before evaluating the
energy. Choose a spanning tree in each component, round the chord flows
to dyadics, subtract their incidence contributions from $b$, and recover
tree flows by leaf elimination. At a leaf the unique remaining tree
flow is its residual nomination, with its orientation sign; transfer
that residual to its parent and delete the leaf. The last residual
vanishes because the component sum of $b$ is zero. This proves exact
conservation. If a conserved reference flow was approximated on every
chord to accuracy $\rho$, each tree correction differs from that
reference by at most $m\rho$: removing a tree edge expresses its flow
as a cut sum, with at most $m$ chord errors. Thus the repair preserves
convergence of increasingly accurate approximations, and it runs in time
polynomial in the input and rounding precision. The numerical helper
implements only the connected simple loopless case and rejects other
graphs; this implementation restriction does not narrow the mathematical
certificate.

For a nonnegative rational $a/b$, a dyadic upper enclosure of its
$k$th root, $k\in\{2,3\}$, is $j/2^q$, where $j$ is the least
nonnegative integer satisfying
\begin{equation}
 j^k b\ge a2^{kq}.
 \label{eq:a-cert-root}
\end{equation}
An integer square root, or binary search for an integer cube root,
followed by one exact upward comparison computes $j$ in polynomial
time in the expanded input length and $q$. Its excess over the root
is less than $2^{-q}$. This constructs the conjugate enclosures and
$\eta$ without algebraic-number arithmetic; potentials can be rounded
independently. These constructions certify achieved accuracy; they give
no polynomial stopping bound for an arbitrary numerical solver.

\subsection{Bregman intervals, pressure objectives, and sign recovery}

\begin{proposition}[Separate-edge rational intervals]
\label{prop:a-cert-bregman}
Assume the certificate in \cref{prop:a-cert-energy} has been verified.
If $\delta>0$, rational endpoints $l_e\le y_e\le u_e$ satisfying
\begin{equation}
 D_e(y_e,l_e)\ge\delta,\qquad D_e(y_e,u_e)\ge\delta
 \label{eq:a-cert-bregman-test}
\end{equation}
certify $x_e^*\in[l_e,u_e]$. If $\delta=0$, $x^*=y$ and the
singleton intervals suffice. Starting from $y_e\pm\eta$, rational
bisection constructs the endpoints in \eqref{eq:a-cert-bregman-test}
to any specified additive boundary accuracy in polynomial bit time.
\end{proposition}
\begin{proof}
Conservation and stationarity give the exact identity
\begin{equation}
 E(y)-E(x^*)=\sum_e D_e(y_e,x_e^*)\le\delta.
 \label{eq:a-cert-bregman-sum}
\end{equation}
Every summand is nonnegative. For fixed $y_e$, differentiation with
respect to $z$ gives
\[
 \frac{d}{dz}D_e(y_e,z)=2c_e^{\sgn(z)}|z|(z-y_e).
\]
The function is continuously differentiable at zero as well. Its
derivative is strictly negative on $z<y_e$ except possibly at the
single point zero, and strictly positive on $z>y_e$ with the same
exception. Consequently it is strictly decreasing to zero at $y_e$
and strictly increasing thereafter. It tends to infinity at both
ends by \eqref{eq:a-cert-scalar-modulus}. Thus its sublevel set at
$\delta$ is one interval. The endpoint inequalities, in the stated
outward directions, enclose this sublevel set and hence $x_e^*$.
The initial endpoints $y_e\pm\eta$ are outside or on the sublevel
boundary because $\beta_L\eta^3/6\ge\delta$. On each side keep
$y_e$ as an inside endpoint and the current outer endpoint with
$D_e\ge\delta$. Bisect and replace the appropriate endpoint using
exact cubic evaluation. After $k$ steps the boundary bracket length
is $\eta2^{-k}$. All rational sizes grow polynomially in $k$.
A zero gap gives $E(y)=E(x^*)$, so uniqueness proves the singleton
claim without a curvature assumption.
\end{proof}

Increasing boundary accuracy at a fixed positive gap does not make
these intervals arbitrarily narrow: it approaches the two roots of
$D_e(y_e,z)=\delta$. Away from zero the divergence grows quadratically
in displacement; near zero its cubic behavior persists. This is the
source of its improvement over a common cube-root radius, and also
its limitation.

The physical edge drop lies in the rational interval
$[G_e(l_e),G_e(u_e)]$. Along a fixed path, sum these intervals with
orientation signs to enclose the terminal potential difference.
For any rational zero-sum weighted potential objective $a^\top\pi$
on a connected graph, choose a spanning-tree flow $q$ with $Aq=a$.
Then $a^\top\pi=q^\top G(x^*)$; signed interval arithmetic on the
tree edges gives its enclosure. Componentwise zero sums give the same
statement for disconnected graphs. These are deterministic-state
certificates on arbitrary graphs. An uncertainty interpretation for
that objective still requires its own graph-dependent reduction.
In particular, an envelope target-drop interval concerns the envelope
law, whose target coefficient may differ from the selected original
resistance scenario.

\begin{proposition}[Posterior endpoint-scenario recovery]
\label{prop:a-cert-signs}
Suppose \eqref{eq:a-cert-energy} is a correctly constructed envelope
for an original target-flow extremum, with allowed resistance endpoints
$\underline\beta_e,\overline\beta_e$ and fixed nominations.
Select $\widehat\beta_e=c_e^+$ if $y_e\ge0$ and
$\widehat\beta_e=c_e^-$ otherwise. Let $[l_e,u_e]$ be verified
intervals containing $y_e$ and $x_e^*$, and put
\[
 a_e=\begin{cases}\max\{0,-l_e\}&y_e\ge0,\\
                       \max\{0,u_e\}&y_e<0,
       \end{cases}
 \qquad r_e=(\overline\beta_e-\underline\beta_e)a_e^2.
\]
If $\widehat x$ is the physical flow at the selected scenario, then
\begin{equation}
 \|\widehat x-x^*\|_\infty^2
 \le\frac{2}{\beta_L}\sum_e r_e.
 \label{eq:a-cert-sign-error}
\end{equation}
In particular, resolved signs give an exactly optimal rational endpoint
scenario even when the exact flow is irrational. More generally, exact
optimality holds whenever every edge satisfies at least one of
\begin{equation}
 \begin{gathered}
 c_e^+=c_e^-=\widehat\beta_e,\qquad
 l_e\ge0\ \hbox{and}\ \widehat\beta_e=c_e^+,\\
 u_e\le0\ \hbox{and}\ \widehat\beta_e=c_e^-,\qquad
 l_e=u_e=0.
 \end{gathered}
 \label{eq:a-cert-compatible}
\end{equation}
\end{proposition}
\begin{proof}
A disagreement between the selected coefficient and the envelope
coefficient at the true sign can occur only on the opposite side of
zero from $y_e$. In that event $|x_e^*|\le a_e$, whence
$|\widehat\beta_e x_e^*|x_e^*|-G_e(x_e^*)|\le r_e$.
Write $d=\widehat x-x^*$ and $D=\|d\|_\infty$.
The two physical potential gradients both annihilate $d$. Pairing
with $d$ and using the scalar monotonicity inequality gives
\[
 \frac{\beta_L}{2}\sum_e|d_e|^3
 \le\sum_e\bigl(G_e(x_e^*)-
              \widehat\beta_e x_e^*|x_e^*|\bigr)d_e
 \le D\sum_e r_e.
\]
If $D=0$ there is nothing to prove; otherwise division by $D$
proves \eqref{eq:a-cert-sign-error}. Under
\eqref{eq:a-cert-compatible}, the selected law and envelope law
agree at $x^*$ on every edge. The envelope state therefore satisfies
the selected scenario's equations, and uniqueness gives
$\widehat x=x^*$. Non-strict signs suffice because both laws vanish
at zero. The envelope theorem supplies the extremum interpretation.
\end{proof}

An upward rational square root of the right side of
\eqref{eq:a-cert-sign-error} certifies scenario loss. Using only
$|y_e-x_e^*|\le\eta$ recovers the uniform bound from
\cref{thm:a-bound-envelope-algorithm}; in particular
$(1+2m\beta_U/\beta_L)\eta$ is a valid, possibly conservative,
loss bound when $\beta_U=\max_e\overline\beta_e$.
There is no universal finite-precision sign-identification guarantee
here: a zero or very small physical flow may keep an interval across
zero. Exact optimality is a sufficient posterior conclusion when
\eqref{eq:a-cert-compatible} succeeds.

\subsection{Conserved-energy support certificates for linear goals}

Separate-edge intervals omit the coupling $Ax=b$. For a rational
linear goal $w^\top x^*$, retain that coupling in the compact convex set
\[
 S(y)=\{x:Ax=b,\ E(x)\le E(y)\}.
\]
This set contains $x^*$, but its other points need not be physical
states of the fixed law.

\begin{theorem}[Rational support certificate and completeness]
\label{thm:a-cert-support}
Let $Ay=b$. For rational $\lambda>0$, $v\in\Q^V$, set
$s=w-A^\top v$. If rational $R_e\ge0$ satisfy
\begin{equation}
 R_e^2\lambda c_e^{\sgn(s_e)}\ge |s_e|^3,
 \qquad
 U_w=\lambda E(y)+v^\top b+\frac23\sum_e R_e,
 \label{eq:a-cert-support}
\end{equation}
then $w^\top x\le U_w$ for every $x\in S(y)$.
Applying the construction to $-w$ gives the interval
$[-U_{-w},U_w]$ for $w^\top x^*$.
For $\lambda=0$, a finite dual bound exists exactly when
$w=A^\top v$; then $w^\top x=v^\top b$ on $Ax=b$ and no
roots are needed.

If $E(y)>E(x^*)$, the infimum of these rational upper bounds equals
$\max_{x\in S(y)}w^\top x$. Hence rational witnesses can approach
the exact support value arbitrarily closely. As conserved trial
energies decrease to $E(x^*)$, the exact support intervals converge
to $w^\top x^*$ for every $w$.
\end{theorem}
\begin{proof}
The scalar calculation in \cref{prop:a-cert-energy}, with coefficient
$\lambda c_e^{\sgn(s_e)}$, gives
\[
 (\lambda E_e)^*(s_e)=
 \frac23\sqrt{\frac{|s_e|^3}{\lambda c_e^{\sgn(s_e)}}}.
\]
Fenchel's inequality gives $s^\top x\le\lambda E(x)+
\sum_e(\lambda E_e)^*(s_e)$. Add $v^\top Ax=v^\top b$
and use $E(x)\le E(y)$. When $\lambda=0$, the conjugate
of the zero function is finite only at $s=0$, proving the stated
condition. These are exactly the goals constant on the conservation
affine space, because $(\ker A)^\perp=\operatorname{im}A^\top$.
They include cut flows and each bridge flow; rational linear elimination
finds $v$ from rational $A,w$ whenever it exists.

To prove completeness of the unrounded dual, minimize $-w^\top x$ subject
to $Ax=b$, $E(x)\le E(y)$. Its value is finite and attained by
coercivity. If $E(y)>E(x^*)$, the point $x^*$ is strictly feasible
for the energy inequality relative to $Ax=b$. Slater duality
therefore gives equality of the support value and the dual infimum,
and a dual minimizer. Redundant conservation rows may first be
removed, so their presence causes no qualification issue.
If $w$ is nonconstant on $Ax=b$, a finite dual minimizer cannot
have $\lambda=0$ by the preceding conjugate calculation.
For $\lambda>0$ the displayed unrounded objective is continuous in
$(\lambda,v)$, including at $s_e=0$: the two asymmetric terms tend
to zero there. Approximate a dual minimizer by rational data keeping
$\lambda>0$, and then take increasingly accurate upward rational
square roots. The resulting valid rational upper bounds tend to the
support value. Constant goals already have an exact rational witness
at $\lambda=0$.

For convergence, the cubic modulus gives the stronger quantitative
statement
\[
 |z_e-x_e^*|^3
 \le \frac{6\bigl(E(y)-E(x^*)\bigr)}{\beta_L}
 \qquad(z\in S(y),\ e\in E).
\]
Indeed, apply \eqref{eq:a-cert-scalar-modulus} at $x^*$ to the
feasible flow $z$. Conservation cancels the summed linear terms, and
$E(z)\le E(y)$ gives the displayed bound. Thus, if $Ay_k=b$ and
$E(y_k)\to E(x^*)$, the whole set $S(y_k)$ eventually lies in any
fixed positive-radius supremum-norm ball about $x^*$. Continuity of
each fixed linear goal proves convergence of its exact support interval.
This estimate is in terms of energy error and gives no rate in
producer runtime or certificate size.
\end{proof}

The support problem has the cubic second-order-cone lift of
\cref{thm:a-bound-envelope-algorithm}, with a linear objective and
an energy sublevel constraint, so numerical conic solvers can propose
its dual variables. The accompanying support module supplies exact
verification, conversion of supplied rational dual data to root
witnesses, and analytic demonstrations; it does not implement a general
conic search for these dual data. The completeness statement gives no
polynomial certificate-size or producer-runtime bound at a specified
support tolerance, and it does not assert that a support interval
always improves a Bregman interval. In particular,
\eqref{eq:a-cert-bregman-sum} uses physical stationarity at $x^*$;
one may not assume the same separate-edge Bregman constraints for
all $x\in S(y)$.

\subsection{Curvature bounds and the constrained Laplacian}

\begin{theorem}[Conservation-aware quadratic goal certificate]
\label{thm:a-cert-hessian}
Let $y\in\Q^m$ satisfy $Ay=b$.
Suppose $0\le E(y)-E(x^*)\le\delta$ is verified, and rational
intervals $[l_e,u_e]$ contain both $y_e$ and $x_e^*$. Define
\begin{equation}
 h_e=\begin{cases}
 2c_e^+l_e&l_e>0,\\
 -2c_e^-u_e&u_e<0,\\
 0&l_e\le0\le u_e.
 \end{cases}
 \label{eq:a-cert-curvature}
\end{equation}
Let $P=\{e:h_e>0\}$ and $Z=\{e:h_e=0\}$. For any rational
$v$ with $s=w-A^\top v$ satisfying $s_Z=0$, put
\begin{equation}
 C(v)=\sum_{e\in P}\frac{s_e^2}{h_e}.
 \label{eq:a-cert-factor}
\end{equation}
Every rational $r\ge0$ with $r^2\ge2\delta C(v)$ proves
\begin{equation}
 |w^\top(x^*-y)|\le r.
 \label{eq:a-cert-goal-radius}
\end{equation}
If $\delta=0$, every goal is exact regardless of $h$ or $s_Z$.
\end{theorem}
\begin{proof}
The second derivative is $E_e''(t)=2c_e^{\sgn(t)}|t|$, including
its continuous zero value. Formula \eqref{eq:a-cert-curvature} is
its minimum on the interval. Taylor's integral formula along the
segment between $x_e^*$ and $y_e$ gives
$D_e(y_e,x_e^*)\ge h_e(y_e-x_e^*)^2/2$. Therefore
\begin{equation}
 \frac12\sum_e h_e(x_e^*-y_e)^2\le\delta.
 \label{eq:a-cert-quadratic-error}
\end{equation}
Since $A(x^*-y)=0$, the goal difference equals $s^\top(x^*-y)$.
The identity $s_Z=0$ and weighted Cauchy--Schwarz on $P$ bound its square by
$C(v)\sum_e h_e(x_e^*-y_e)^2\le2\delta C(v)$.
Empty sums are zero, so the proof also covers $P=\varnothing$.
For a zero gap use strict convexity, which gives $x^*=y$ directly.
\end{proof}

\begin{proposition}[Optimal factor and zero-curvature obstruction]
\label{prop:a-cert-projection}
Write $D=\operatorname{diag}(1/h_e:e\in P)$. The zero-edge
condition is feasible exactly when
\begin{equation}
 w_Z^\top z=0\quad\hbox{for every }z\in\ker A_Z.
 \label{eq:a-cert-zero-obstruction}
\end{equation}
When feasible, the minimum $C_*$ of \eqref{eq:a-cert-factor} is
attained by a rational solution of
\begin{equation}
 \begin{bmatrix}A_PD A_P^\top&A_Z\\ A_Z^\top&0\end{bmatrix}
 \begin{bmatrix}v\\\xi\end{bmatrix}
 =\begin{bmatrix}A_PD w_P\\w_Z\end{bmatrix}.
 \label{eq:a-cert-kkt}
\end{equation}
Singular potential gauges and redundant zero-edge constraints are
allowed. The value $\sqrt{2\delta C_*}$ is the exact supremum of
$|w^\top d|$ over
\begin{equation}
 Ad=0,\qquad \tfrac12\sum_e h_ed_e^2\le\delta
 \label{eq:a-cert-error-set}
\end{equation}
for every $\delta>0$. If \eqref{eq:a-cert-zero-obstruction} fails,
that supremum is infinite. Thus optimizing the Laplacian constant is
sharp for this quadratic information, without asserting sharpness for
the physical error or for the additional interval constraints.
\end{proposition}
\begin{proof}
The equation $A_Z^\top v=w_Z$ is solvable exactly when $w_Z$
is orthogonal to $\ker A_Z$, by the fundamental subspace identity.
If it is inconsistent there is $z\in\ker A_Z$ with
$w_Z^\top z\ne0$. Extending $z$ by zeros on $P$ gives a
zero-curvature circulation; arbitrary multiples satisfy
\eqref{eq:a-cert-error-set} and have unbounded goal.

For the feasible case, minimize the squared norm
$\|D^{1/2}(w_P-A_P^\top v)\|^2$ on the affine space
$A_Z^\top v=w_Z$. Its affine image in Euclidean space is closed,
so a least-norm point exists. First-order optimality on that affine
space says its gradient lies in $\operatorname{im}A_Z$.
Writing a multiplier $\xi$ gives exactly
\eqref{eq:a-cert-kkt}; conversely these equations imply global
optimality of the convex quadratic. The system is rational and
consistent. Exact elimination with free coordinates set to zero
returns rational $v,\xi$, even with redundant rows. Determinant
bounds after clearing denominators give polynomial output bit length
and polynomial bit complexity; a singular system need not be perturbed.

For sharpness, take an optimal solution and put
$d_P=D(w_P-A_P^\top v)$ and $d_Z=-\xi$.
The first row of \eqref{eq:a-cert-kkt} gives $Ad=0$, and the
second gives $s_Z=0$. Consequently
\[
 w^\top d=s^\top d=C_*,\qquad \sum_e h_ed_e^2=C_*.
\]
If $C_*>0$, multiply $d$ by $\sqrt{2\delta/C_*}$ to attain
the upper bound from weighted Cauchy--Schwarz. If $C_*=0$,
$s=0$ on all edges, so $w=A^\top v$ and every circulation goal
is zero. This also handles $P=\varnothing$: consistency makes
$w$ a cut goal, and the optimal factor is the empty sum zero.
\end{proof}

When every $h_e>0$, let $H=\operatorname{diag}(h)$ and
$L_H=AH^{-1}A^\top$. The minimum can equivalently be written
\begin{equation}
 C_*=w^\top H^{-1}w-
 (AH^{-1}w)^\top L_H^+(AH^{-1}w),
 \label{eq:a-cert-effective}
\end{equation}
where $L_H^+$ is the Moore--Penrose inverse. To verify this expression,
expand $C(v)$ as a quadratic in $v$, solve
$L_Hv=AH^{-1}w$, and substitute. The right side is in
$\operatorname{im}L_H=\operatorname{im}A$, so the substitution
is unaffected by component gauges. For a single-edge selector
$w=\one_e$, writing $a_e$ for that incidence column gives
\[
 C_*=\frac1{h_e}-\frac{R_{\rm eff}(e)}{h_e^2},\qquad
 R_{\rm eff}(e)=a_e^\top L_H^+a_e.
\]
Here the auxiliary network has conductances $1/h_e$.
For a bridge the unit terminal current has no alternative route,
so $R_{\rm eff}(e)=h_e$ and the factor vanishes, as conservation
also proves. This is the usual electrical projection and effective
resistance identity, now used as a goal-error constant.

The supplied producer solves \eqref{eq:a-cert-kkt} by dense exact
rational elimination. It is intended for modest certificate sizes;
no sparse scalability claim follows. Certificate validity does not
require optimal potentials: with positive curvatures, even a rounded
numerical Laplacian solution gives a valid, perhaps larger, factor.
With zero curvatures, all equalities $A_Z^\top v=w_Z$ must instead
hold exactly, and regularizing away an inconsistent constraint cannot
produce a valid witness. The implementation first handles zero energy
gap, then rejects the zero-curvature obstruction. The full pipeline may
omit this optional goal witness and retain its independently verified
Bregman interval.

\subsection{An exact comparison on two long paths}

\begin{example}[Information lost by separate-edge bounds]
\label{ex:a-cert-paths}
Consider two internally disjoint source--sink paths, each of length
$L\ge1$, with every edge oriented toward the sink, all coefficients
one, source nomination two, sink nomination minus two, and zero internal
nominations. The physical flow is one on every edge: its path drops
are both $L$, and strict convexity gives uniqueness. For rational
$0<\eps<1$, choose $y=1+\eps$ on the first path and
$y=1-\eps$ on the second. Its exact energy gap is
\begin{equation}
 \delta=\frac L3\bigl((1+\eps)^3+(1-\eps)^3-2\bigr)
       =2L\eps^2.
 \label{eq:a-cert-path-gap}
\end{equation}
Integral physical potentials, decreasing by one on each path edge,
give this exact dual energy value, with rational conjugate roots.

Every conserved flow has path values $t,2-t$. Its energy is strictly
convex, symmetric about $t=1$, and strictly increasing with $|t-1|$.
Consequently $S(y)$ is exactly $1-\eps\le t\le1+\eps$,
even though the energy formula uses absolute cubes outside $[0,2]$.
The exact support interval of a first-path edge is therefore of width
$2\eps$. Its upper dual can also be attained rationally. Choose
\[
 \lambda=\frac1{4L\eps},\qquad
 (A^\top v)_e=w_e-\lambda t_e^2,
\]
where $t_e=1+\eps$ on the first path and $1-\eps$ on the second,
and $w$ selects one first-path edge. The two sums of these prescribed
potential drops agree because
$1-\lambda L(1+\eps)^2=-\lambda L(1-\eps)^2$.
Thus rational node potentials exist. The residual is
$s_e=\lambda t_e^2$, with exact root
$R_e=\lambda t_e^3$; equality in Fenchel and in the energy
constraint gives $U_w=1+\eps$. Exchanging the path values and
using $-w$ gives the lower support bound as well.

For fixed $L$ and rational $\eps\downarrow0$, let $z=y_e+a$ remain
positive. Direct expansion gives
\[
 D_e(y_e,y_e+a)=y_ea^2+\frac23a^3.
\]
The two boundary displacements satisfying $D_e=2L\eps^2$
therefore obey $a/\eps\to\pm\sqrt{2L}$. The separate-edge
Bregman width is
\begin{equation}
 2\sqrt{2L}\,\eps+o(\eps).
 \label{eq:a-cert-path-bregman}
\end{equation}
All interval endpoints tend to one, so their curvatures tend to two.
The circulation space is one-dimensional, generated by $d$ equal to
$+1$ on one path and $-1$ on the other. Directly optimizing
\eqref{eq:a-cert-error-set}, or using
\cref{prop:a-cert-projection}, gives
$C_*=1/\sum_e h_e\to1/(4L)$ for the selected edge.
Thus the optimized Laplacian width is
\begin{equation}
 2\sqrt{2\delta C_*}=2\eps+o(\eps).
 \label{eq:a-cert-path-hessian}
\end{equation}
These asymptotics hold for exact boundaries. They also hold for rational
outer endpoints and roots with errors $o(\eps)$.
They are pointwise in $L$. Choosing $L$ large and then $\eps$
sufficiently small makes the improvement factor arbitrarily large;
no uniform-in-$L$ asymptotic or computational lower bound is implied.
The physical solution of this family is elementary.
\end{example}

The centers of these intervals differ. The Laplacian interval is
centered at $w^\top y$, whereas the exact support interval need not
be. Equal asymptotic widths therefore do not imply equal endpoints.
For a concrete warning against combining their hypotheses, take
$L=1$ and $\eps=1/10$. The opposite support endpoint has
$D_1(11/10,9/10)=29/750>\delta=1/50$. It belongs to $S(y)$
but fails a separate-edge physical Bregman constraint. Intersecting
independently verified enclosures of $x^*$ is valid; silently
restricting the support problem using those physical constraints is
a different optimization problem.

\subsection{Verification of the original uncertainty instance}

The preceding certificates validate a deterministic energy model. Its
JSON data alone cannot establish that it is the right envelope for a
specified uncertainty problem. The complete pipeline additionally
verifies a connected simple loopless graph with no $K_4$ minor,
fixed balanced nominations, independent positive rational resistance
intervals or explicit finite sets, a target edge, and maximum or
minimum direction. It accepts no operating side constraints or
correlated uncertainty. Its simple-graph restriction is narrower than
the multigraph scope of the mathematics.

The verifier checks indices, simplicity, connectivity, dimensions, and
exact balance. Allowed sets must be nonempty and strictly positive. It recognizes
its graph class by repeatedly eliminating a vertex of degree at most
two and joining its two neighbors when needed. The $K_4$-minor-free
characterization underlying \cref{lem:a-bound-signs} implies that
such a vertex always exists in the remaining graph. Suppressing a
degree-two vertex is a minor operation; degrees zero or one also
preserve the class. Conversely a completed elimination order reverses
to adding vertices whose existing neighbors form cliques of size at
most two. Starting with bags of at most three vertices and attaching
a new bag at its neighbor clique gives a tree decomposition of width
at most two; therefore the original graph has no $K_4$ minor.
This also explains why the test accepts dangling cyclic blocks.

Next solve the grounded all-unit-conductance Laplacian for a unit
injection at the target tail and withdrawal at its head. Exact
rational elimination computes every adjoint drop and its sign,
including exact zeros. Let these signs be $\sigma_e$, and let
$d=1$ for maximization and $d=-1$ for minimization. Set
\[
 t_e=d\sigma_e\quad(e\ne a),\qquad t_a=-d.
\]
The directional coefficients are
\[
 c_e^+=\begin{cases}\overline\beta_e&t_e>0,\\
                         \underline\beta_e&t_e\le0,\end{cases}
 \quad
 c_e^-=\begin{cases}\overline\beta_e&t_e<0,\\
                         \underline\beta_e&t_e\ge0.\end{cases}
\]
For finite sets, $\underline\beta_e,\overline\beta_e$ are their
minimum and maximum. The target's reversed rule and the zero-sign
choice are those of \cref{thm:a-bound-envelope-structural}.
The confluence and finite-secant proofs there justify this exact
mapping for all positive secant resistances. The verifier recomputes
it from the original instance, rather than trusting stored coefficients
or floating-point signs.

It then checks an energy witness for this envelope and a second energy
witness for the selected original endpoint scenario. Each yields
Bregman intervals, optionally intersected with independently verified
linear-goal intervals for the reconstructed target selector. Empty
intersections are rejected. Let the resulting target intervals be
$I^*=[l^*,u^*]$ and $I^\beta=[l^\beta,u^\beta]$.
The original-scenario loss is bounded directly by
\begin{equation}
 \begin{cases}
 0\le\OPT-x_a(\beta)\le u^*-l^\beta&\text{for maximization},\\
 0\le x_a(\beta)-\OPT\le u^\beta-l^*&\text{for minimization}.
 \end{cases}
 \label{eq:a-cert-direct-loss}
\end{equation}
Indeed the original envelope theorem identifies $\OPT$ with $x_a^*$,
the scenario belongs to the allowed product, and interval subtraction
bounds their difference. If \eqref{eq:a-cert-compatible} holds, the
checker instead reports zero loss by exact endpoint optimality.
A zero loss certifies the selected scenario, not an exact numerical
value for its possibly irrational flow.

The second witness requires no second physical-flow optimization.
The producer may reuse the conserved envelope trial flow, round it
again if necessary, and recompute a dual potential proposal under the
selected symmetric law. Its independent primal-dual check certifies
the resulting scenario interval; poor reuse only widens it. This is a
posterior alternative to the uniform recovery factor.

\subsection{Scaling, implementation scope, and reproducible evidence}

Let $B=\sum_{b_v>0}b_v$ and $C=\max_{e,\sigma}c_e^\sigma$.
For $B>0$ the producer solves and certifies normalized data
$\widetilde b=b/B$, $\widetilde c=c/C$. The exact lift is
\begin{equation}
 \begin{gathered}
 y=B\widetilde y,\quad p=CB^2\widetilde p,\quad
 \zeta=CB^3\widetilde\zeta,\quad
 \delta=CB^3\widetilde\delta,\quad\eta=B\widetilde\eta,\\
 [l_e,u_e]=B[\widetilde l_e,\widetilde u_e],\quad
 h_e=CB\widetilde h_e,\quad
 v=\widetilde v,\quad C(v)=\widetilde C(v)/(CB),\quad
 r=B\widetilde r.
 \end{gathered}
 \label{eq:a-cert-scaling}
\end{equation}
Here $\zeta$ denotes the vector of conjugate-root bounds, and
$l_e,u_e$ are interval endpoints. Cubic homogeneity
proves the energy and root identities: $|A^\top p|^3/c$ scales
by $C^2B^6$. The second-derivative formula proves the curvature
identity, and $r^2\ge2\delta C(v)$ then scales by $B^2$.
Thus every rational inequality is preserved. Certifying in normalized
units before this exact lift makes the rounding error covariant under
common rational changes of flow and resistance units. Normalizing
only the floating-point solve would not have that property.
When $B=0$, balance implies $b=0$ and the producer uses the exact
zero physical state. A positive-gap zero-load trial circulation is
still a legitimate input to the general verifier and may exhibit the
zero-curvature obstruction.

The numerical energy producer combines the explicit cubic cone lift
(solved with CVXPY and Clarabel), conservation repair, rounded
least-squares potential proposals, and exact root bounds. It is a
practical proposal mechanism, not an implementation of the rational
convex-body algorithm of \cref{thm:a-bound-envelope-algorithm}, and it
offers no guarantee for a requested accuracy. Extreme coefficient
ratios can cause large certified intervals or a failed numerical solve.
Verification uses only Python's standard library and explicit
exceptions, so its mathematical checks remain active under
\texttt{python -O}. The original-instance reader also rejects duplicate
JSON keys, unknown fields, nonstring rational fields, and Boolean
indices. The simpler base-energy reader certifies the deterministic
instance encoded in its own file and does not enforce this complete
original-instance schema. The Bregman helper requires prior base
verification, and the integrated goal verifier repeats that check.

The implementation files, in \texttt{code/potential\_flow\_mpd/},
have distinct roles:
\begin{center}
\begin{tabular}{@{}p{.45\linewidth}p{.48\linewidth}@{}}
\toprule
File & Verified role\\
\midrule
\texttt{envelope\_rational\_\allowbreak certificates.py} &
Deterministic energy, conjugate roots, conserved flow, and uniform radius.\\
\texttt{envelope\_bregman\_bounds.py} &
Outward scalar interval sharpening after base verification.\\
\texttt{goal\_flow\_certificate.py} &
Support verification and exact constrained-Laplacian production and verification.\\
\texttt{certified\_envelope.py} &
Original data, exact envelope mapping, two state certificates,
optional target witnesses, and scenario-loss interpretation.\\
\bottomrule
\end{tabular}
\end{center}

The saved six-edge deterministic witness gives a uniform interval width
below $3.405\cdot10^{-4}$. Rechecking it and sharpening by exact
Bregman arithmetic gives maximum edge width below
$2.894\cdot10^{-6}$ and maximum edge-drop width below
$4.0\cdot10^{-6}$; all six edge intervals exclude zero.
Applying the exact Laplacian producer without another physical solve
gives maximum goal width below $9.447\cdot10^{-7}$.
These are widths of verified enclosures, not measurements of the
true numerical error. Exact optimality of an original scenario would
additionally require an independently checked envelope mapping for this
saved base witness.

The two-path demonstration at $\eps=10^{-3}$ gives the following
rounded displays of exactly checked bounds. The support column is the
analytic width from \cref{ex:a-cert-paths}.
\begin{center}
\begin{tabular}{@{}rrrr@{}}
\toprule
$L$ & Separate-edge width & Laplacian width & Exact support width\\
\midrule
2 & $0.00399800593$ & $0.00200200435$ & $0.002$\\
8 & $0.00799603843$ & $0.00200401743$ & $0.002$\\
32 & $0.0159922895$ & $0.00200807006$ & $0.002$\\
\bottomrule
\end{tabular}
\end{center}

A replay of the original-instance benchmark produced valid certificates
in ten cases, including maximum and minimum directions, mixed finite
and interval sets, reversed orientations, an 80-edge graph of rank
39, a full coefficient ratio of $3\cdot10^6$, supply $10^{-8}$, a dangling
zero-adjoint block, zero nominations, and an ambiguous-sign case.
Nine cases certified exact endpoint optimality. Some numerical solves
reported \texttt{optimal\_inaccurate}; the reported guarantees came
from the exact witnesses. The ambiguous-sign instance retained a
nonzero certified loss: the goal witness reduced its target enclosure
width from about $4.99\cdot10^{-4}$ to $4.37\cdot10^{-6}$
and its loss upper bound from about $6.92\cdot10^{-4}$ to
$5.28\cdot10^{-6}$, approximately $114$-fold and $131$-fold reductions
in verified bounds. This says nothing about the actual scenario error
or exact sign resolution.

One case adapts the topology and pipe geometry of the open water-network
dataset accompanying \citet{vrachimis2018-hydraulic}.
It uses 11 nodes and 13 pipes, synthetic unit demands at the nine
nonreservoir nodes numbered 2 through 10, zero demand at the leak
subdivision node, and supply nine at node 1. Its quadratic proxy
coefficient is $(\mathrm{length}/1000)(300/\mathrm{diameter})^5$,
with a $\pm10\%$ interval. The target is original pipe 11,
from node 4 to node 9. This transformed instance tests an external
topology; it does not reproduce the source's hydraulic law,
measurements, demand scenarios, or leakage estimates, and it
establishes neither hydraulic prediction accuracy nor large-network
scalability.

The retained verification record includes ordinary and optimized
Python runs of independently written deterministic, Bregman, goal,
and original-instance checks. Their controls include exact signed
states, zero gaps, zero-curvature cut and cycle goals, parallel edges,
disconnected gauges, malformed witnesses, and all 771 connected
labeled graphs on two through five vertices tested against the
branch-set definition of a $K_4$ minor. The two small numerical cases
were additionally compared with all 128 endpoint scenarios combined,
using a separate numerical physical solver; this latter comparison is
a regression check, not a proof. Exact scale tests at $10^{400}$
and $10^{-400}$ verify the covariance in
\eqref{eq:a-cert-scaling}. A further diagnostic checks the sharp
quadratic constant directly in circulation coordinates, including
singular zero-edge cases, and the rational two-path support equalities.
These finite checks supplement the proofs and do not establish their
universal statements.

From the repository root, the saved complete certificate can be checked
without scientific packages by
\begin{quote}\small\ttfamily
cd code/potential\_flow\_mpd\\
python -S -O certified\_envelope.py verify \textbackslash\\
\hspace*{1em}certified\_envelope\_example.json
\end{quote}
The deterministic saved witness is checked with the base module's
\texttt{--verify} option; the goal module runs the analytic demonstration
with no arguments and accepts \texttt{--certificate} for the saved
six-edge comparison. The numerical benchmark entry point is
\texttt{certified\_envelope\_benchmarks.py}; it requires the scientific
Python environment and writes example and result JSON files beside
itself. The retained replay uses an isolated copy so the distributed
saved witnesses remain unchanged. Source hashes, full commands,
interpreter and package versions, exact rational comparison data,
and return codes accompany the verification record. No solver status,
rounded display, or historical benchmark result substitutes for
rechecking the exact witness against its original instance.
